\subsection{A full elementary inertia group for every nonordinary principal surface}
\label{sec:nonordinary-uniform-Q-inertia-candidate}
\status{Proved and independently checked.}

Throughout this subsection the ground field is algebraically closed of
characteristic two. Use the nine theta coordinates and the four direction
coefficients $\alpha_i$ of Lemma~\ref{lem:DE-A3-theta-pfaffian}. Write
$M(s)$ for the alternating contraction matrix, $q=q_\alpha$ for its
Pfaffian cubic, and $Q(s)=\sum_i s_i^3$. The actual nonordinary principal
surface tensor has $q\ne0$, by
Proposition~\ref{prop:DE-A3-ppav-degree-one}. The relative Frobenius
bundle is the Euler bundle $\Omega^1_{\mathbf P^8}$ and its determinant
is $q^2Q$, by \eqref{eq:DE-A3-relative-theta-Hasse}.

\begin{lemma}[Stable loss one on the Hermitian cubic]
\label{lem:nonordinary-uniform-Q-stable-loss-candidate}
For every four-parameter tensor with $q\ne0$, the relative Frobenius
has linear rank and stable rank seven at the generic point of $Q=0$.
Its determinant has order one there, and the linearized kernel map
from that divisor to the constant projective eight-space is $[s]\mapsto[s]$.
\end{lemma}
\begin{proof}
The cubic $Q$ is smooth, since its partial derivatives are $s_i^2$;
it is therefore geometrically integral. Also $q$ is not proportional
to $Q$. Indeed, if all its line coefficients vanished while its pure
cube coefficient were nonzero, all four $\alpha_i$ would be nonzero
and would satisfy $\sum_{j\ne i}\alpha_j^3=0$. Subtracting any two
of these equations makes their cubes equal, but then each sum of the
other three equals that common nonzero cube, a contradiction.
Consequently $Q\setminus\{q=0\}$ is dense and the determinant has
order one along $Q$.

At such a point $\ker M=ks$ and $\operatorname{im}M=s^\perp$.
Since $Q=0$, there is a vector $u$ with $Mu=s^{[2]}$. On the dual
relative Cartier quotient its linearized kernel is $k\bar u$ in
$W/ks$. The rank-seven first image is already stable precisely when
\begin{equation}\label{eq:nonordinary-Q-stable-test-candidate}
                          s^t u^{[2]}\ne0.
\end{equation}
To check this criterion directly, the kernel of the next squared
linearization is represented by $u^{[2]}$ modulo $ks^{[2]}$.
The latter line already belongs to $\operatorname{im}M$.
Thus this next kernel meets the first image if and only if
$u^{[2]}\in\operatorname{im}M=s^\perp$. Equality of the first two
linearized image ranks implies stabilization of all subsequent ranks.
Changing $u$ by a multiple of $s$ does not change the test, because
$Q=0$.

First suppose that for some direction, renamed vertical, its line
coefficient
\[
          c_0=\alpha_0(\alpha_1^3+\alpha_2^3+\alpha_3^3)
\]
is nonzero and two of the other cubes differ. The stabilizer of a
direction in $\GL_2(\mathbf F_3)$ permutes its other three directions
as the full symmetric group. We may therefore arrange
$\alpha_1^3+\alpha_3^3\ne0$. Order coordinates lexicographically
by $(x,y)\in\mathbf F_3^2$, with the vertical direction first, followed
by slopes zero, one, two. Take the pencil
\[
                 s=(1,1,z,z,0,0,0,0,0).
\]
Here $Q=0$ identically and $q=c_0z$. Delete row and column zero of
$M$ to obtain the invertible alternating eight-by-eight matrix $B$.
Its Pfaffian is $q$. The inverse alternating pairing between
$e_0+e_1$ and $e_2+e_3$, with coordinate zero deleted, has numerator
\[
 P_6(0,1,2)+P_6(0,1,3)
                    =(\alpha_1^3+\alpha_3^3)z,
\]
where $P_6(i,j,k)$ denotes the Pfaffian after deleting those three
indices. For completeness the two Pfaffians are
\[
 \begin{split}
 P_6(0,1,2)&=(\alpha_1^3+\alpha_2^3+\alpha_3^3)z
                         +\alpha_1\alpha_2\alpha_3z^3,\\
 P_6(0,1,3)&=\alpha_2^3z+\alpha_1\alpha_2\alpha_3z^3.
 \end{split}
\]
These identities follow by the fifteen matching terms for each
six-by-six Pfaffian. The inverse-matrix Pfaffian formula, or its
entrywise product with $B$, now gives a solution $u_0=0$ of
$Mu=s^{[2]}$ satisfying
\[
 s^t u^{[2]}=(z^4+z)
       \left(\frac{(\alpha_1^3+\alpha_3^3)z}{c_0z}\right)^2.
\]
One can also derive this last formula without expanding the solution:
write $s=v+zw$ with binary vectors $v=e_0+e_1$, $w=e_2+e_3$.
The inverse of $B$ is alternating; its self-pairings on $v,w$ vanish,
while its two cross-pairings agree. Expanding
$s^t(B^{-1}s^{[2]})^{[2]}$ gives exactly the displayed expression.
It is nonzero in $k(z)$, proving the required generic stable rank.

Here is the complete coefficient reduction. Put $x_i=\alpha_i^3$.
If exactly two $x_i$ are nonzero, either nonzero direction has
$c_i\ne0$ and two unequal remaining cubes. If exactly three are
nonzero and not equal, at least one nonzero $x_i$ differs from
$S=\sum_jx_j$; that direction has $c_i\ne0$, and its remaining
cubes include zero and a nonzero value. Three equal nonzero cubes,
or just one nonzero cube, are precisely excluded cases $q=0$.
If all four are nonzero, some $x_i\ne S$, since otherwise $S=4S=0$.
Choose such an $i$. If its other three cubes are not equal we are
done. If they equal $v$ and $x_i\ne v$, choose one of these three
instead: its complementary sum is $x_i\ne0$ and its complementary
cubes are unequal. The only remaining case is therefore
\[
                         x_0=x_1=x_2=x_3\ne0.
\]

In this last case a scalar multiple of the tensor and a unitary
diagonal theta change reduce all coefficients to one. To see this
explicitly, divide by one coefficient and write the others as powers
of a primitive cube root $\zeta$. A diagonal change
$e_a\mapsto\zeta^{R(a)}e_a$, with $R$ a binary quadratic form over
$\mathbf F_3$, multiplies the direction coefficient by
$\zeta^{-R(v_i)}$. The evaluations at the four directions form exactly
the hyperplane of four-tuples with sum zero: the formula for
$R=Ax^2+Bxy+Cy^2$ gives the four values
$C,A,A+B+C,A-B+C$. Adding a common exponent supplies its complementary
one-dimensional space. Thus such a diagonal change, together with a
common tensor scalar, makes every coefficient one. The diagonal change
satisfies $g^{[2]}=g^{-t}$, so preserves $Q$ and the typed relative
Frobenius system. Multiplying its Frobenius matrix by a nonzero
constant preserves every iterated rank.

For this single normal form an exact point over $\mathbf F_8$ suffices.
Let $\theta^3+\theta+1=0$, and encode
$a_0+a_1\theta+a_2\theta^2$ by $a_0+2a_1+4a_2$.
Take
\[
 \begin{split}
 s&=(0,5,3,3,6,0,4,7,5),\\
 u&=(3,3,3,1,5,7,6,0,0).
 \end{split}
\]
The contraction rule gives, in that field,
\[
        Q(s)=0,\qquad q(s)=7,\qquad Mu=s^{[2]},
                         \qquad s^tu^{[2]}=7.
\]
These are exact finite-field equalities, not a sampling inference.
They establish \eqref{eq:nonordinary-Q-stable-test-candidate} in the
remaining normal form. Openness of the second-iterate rank-seven
condition on the integral $Q$ proves the general assertion.
Finally on $Q\setminus\{q=0\}$ the kernel of the coherent
linearization, whose domain is $\operatorname{ann}(s^{[2]})$,
is exactly $ks$. Its projective map is the stated identity.
\end{proof}

\begin{corollary}[Full seven-dimensional elementary geometric inertia]
\label{cor:nonordinary-ppav-uniform-root-inertia-candidate}
For every nonordinary principally polarized abelian surface in
characteristic two, the degree-one relative geometric image on the
maximal-rank smooth-member open of $|3\Theta|$ contains the entire
group of order $2^7$ fixing a seven-dimensional hyperplane pointwise
and acting trivially on its one-dimensional quotient.
\end{corollary}
\begin{proof}
Apply Lemma~\ref{lem:moving-kernel-full-root-inertia-candidate}
to $Q\subset\mathbf P^8$ and the relative Euler bundle, which embeds
in a constant nine-dimensional bundle on an affine coefficient chart.
The preceding lemma verifies stable rank seven, a one-dimensional
linearized kernel, and determinant order one. Its kernel map is the
identity on the smooth seven-dimensional cubic $Q$, so its differential
has rank seven. The lemma gives all $2^7$ elementary inertia elements.
Restricting the coefficient invertibility open to its dense smooth-curve
open preserves the finite geometric image. All assertions are geometric
and therefore remain valid after algebraically closed extension.
\end{proof}

\begin{remark}[Certificate and remaining global gate]
\label{rem:nonordinary-ppav-uniform-root-scope-candidate}
The exact symbolic certificate is
\path{agents/monodromy/code/ppav_theta_Q_pencil_pfaffians_complete.py};
it verifies the two six-Pfaffians, every coordinate of $Mu=q s^{[2]}$
for the numerator vector, and its pairing, in
$\mathbf F_2[\alpha_0,\ldots,\alpha_3,z]$.
The exact normal-form witness is independently reconstructed by
\path{agents/monodromy/code/ppav_theta_equal_cubes_F8_witness.py}.
Their initial author executions took respectively
$0.1280531329975929$ and $0.0005276640004012734$ seconds.
The earlier bounded finite-field exploration and both failed binary
isotropic-pencil searches are retained separately; neither is a proof
input. This corollary does not assert full monodromy: connectivity of
the nonzero Lang cover, or another global irreducibility argument,
remains to be proved for an arbitrary fixed nonproduct principal surface.
\end{remark}
